% app_g 附录J 非坐标基 (书页 483-494 / p498-p509)
%--F-TAIL--
% 附录J Noncoordinate Bases
\chapter{非坐标基}

% ---- 书页 483 / p498 ----
在研究流形的早期，我们曾决定为切空间选取与坐标相适配的基。既出于审美也出于实用的理由，我们应当再一次考察联络与曲率的形式体系，但这一次在切空间中使用一组\emph{并非}取自任何坐标系的基矢量。将会看到，这一侧重点上的小小改变，会揭示出看待联络与曲率的另一种视角；在这种视角下，它与粒子物理中规范理论（gauge theories）的联系要透明得多。事实上，将要引入的概念非常直截了当，只是这套记号是一场噩梦，所以它看起来比实际上更难。

迄今为止，我们一直在利用这样一个事实：在点 $p$ 处，切空间 $T_{p}$ 的一个自然基由对该点坐标的偏导数给出，即 $\hat{e}_{(\mu)} = \partial_{\mu}$；类似地，余切空间 $T_{p}^{*}$ 的一个基由坐标函数的梯度给出，即 $\hat{\theta}^{(\mu)} = \mathrm{d}x^{\mu}$。然而，没有什么能阻止我们随自己的心意搭建任意的基。因此，让我们设想在流形的每一点引入一组基矢量 $\hat{e}_{(a)}$（用拉丁字母而非希腊字母作指标，以提醒我们它们不与任何坐标系相联系）。我们将把这些基矢量选成“正交归一的”（orthonormal），其含义与我们所工作的流形的号差相适配。也就是说，如果度规的规范形式写作 $\eta_{ab}$，我们就要求基矢量的内积为
\begin{equation*}
g(\hat{e}_{(a)}, \hat{e}_{(b)}) = \eta_{ab},
\tag{J.1}
\end{equation*}
其中 $g(\ ,\ )$ 是通常的度规张量。于是，在洛伦兹型时空中 $\eta_{ab}$ 代表 Minkowski 度规，而在具有正定度规的空间中它则代表欧几里得度规。构成正交归一基的这组矢量有时被称为一个 \textbf{tetrad}（即“四标架”，源自希腊语 tetras，“四个一组”）或 \textbf{vielbein}（源自德语，意为“许多条腿”）。在不同的维数下，它偶尔也被称为 \emph{vierbein}（四）、\emph{dreibein}（三）、\emph{zweibein}（二），等等。正如我们一般找不到覆盖整个流形的坐标卡片，我们也常常无法找到一组处处有定义的光滑基矢量场。照例，我们可以通过在不同的片（patch）上工作、并确保各片在交叠处表现良好来克服这个问题。

拥有一个基的意义在于：任何矢量都可以表示为基矢量的线性组合。具体地说，我们可以把原来的基矢量

% ---- 书页 484 / p499 ----
$\hat{e}_{(\mu)} = \partial_{\mu}$ 用新的基表示出来：
\begin{equation*}
\hat{e}_{(\mu)} = e_{\mu}{}^{a}\hat{e}_{(a)}.
\tag{J.2}
\end{equation*}
分量 $e_{\mu}{}^{a}$ 构成一个 $n\times n$ 的可逆矩阵。（按照我们通常把对象与其分量混为一谈的做法，我们将把 $e_{\mu}{}^{a}$ 称为 tetrad 或 vielbein，而且常用其复数形式“vielbeins”。）我们把指标换位来表示它们的逆，得到 $e^{\mu}{}_{a}$，它们满足
\begin{equation*}
e^{\mu}{}_{a}e_{\nu}{}^{a} = \delta^{\mu}_{\nu}, \qquad e_{\mu}{}^{a}e^{\mu}{}_{b} = \delta^{a}_{b}.
\tag{J.3}
\end{equation*}
它们正是矢量 $\hat{e}_{(a)}$ 在坐标基下的分量：
\begin{equation*}
\hat{e}_{(a)} = e^{\mu}{}_{a}\hat{e}_{(\mu)}.
\tag{J.4}
\end{equation*}
用逆 vielbein 表示，式 (J.1) 变为
\begin{equation*}
g_{\mu\nu}e^{\mu}{}_{a}e^{\nu}{}_{b} = \eta_{ab},
\tag{J.5}
\end{equation*}
或者等价地
\begin{equation*}
g_{\mu\nu} = e_{\mu}{}^{a}e_{\nu}{}^{b}\eta_{ab}.
\tag{J.6}
\end{equation*}
这最后一个方程有时使人们说，vielbein 是度规的“平方根”。

我们可以类似地在 $T_{p}$ 中搭建一组正交归一的 1-形式基，记作 $\hat{\theta}^{(a)}$。可以选取它们与基矢量相容，也就是说
\begin{equation*}
\hat{\theta}^{(a)}(\hat{e}_{(b)}) = \delta^{a}_{b}.
\tag{J.7}
\end{equation*}
一个直接的推论是：这组正交归一 1-形式与其基于坐标的表亲 $\hat{\theta}^{(\mu)} = \mathrm{d}x^{\mu}$ 之间的关系为
\begin{equation*}
\hat{\theta}^{(\mu)} = e^{\mu}{}_{a}\hat{\theta}^{(a)}
\tag{J.8}
\end{equation*}
以及
\begin{equation*}
\hat{\theta}^{(a)} = e_{\mu}{}^{a}\hat{\theta}^{(\mu)}.
\tag{J.9}
\end{equation*}
于是，vielbein $e_{\mu}{}^{a}$ 身兼二职：既是坐标基矢量用正交归一基矢量表出的分量，又是正交归一基 1-形式用坐标基 1-形式表出的分量；而逆 vielbein 则既是正交归一基矢量用坐标基表出的分量，又是坐标基 1-形式用正交归一基表出的分量。

% ---- 书页 485 / p500 ----
任何其他矢量也可以用它在正交归一基中的分量来表示。如果矢量 $V$ 在坐标基中写作 $V^{\mu}\hat{e}_{(\mu)}$，在正交归一基中写作 $V^{a}\hat{e}_{(a)}$，那么这两组分量之间的关系为
\begin{equation*}
V^{a} = e_{\mu}{}^{a}V^{\mu}.
\tag{J.10}
\end{equation*}
这样，vielbein 让我们可以“在拉丁指标和希腊指标之间来回切换”。张量的一个美妙性质——基于指标的位置，通常只有一种合理的做法——在这里帮了大忙。我们可以进一步用两种基中的任何一种来标记多指标张量，甚至用混合分量来标记：
\begin{equation*}
V^{a}{}_{b} = e_{\mu}{}^{a}V^{\mu}{}_{b} = e^{\nu}{}_{b}V^{a}{}_{\nu} = e_{\mu}{}^{a}e^{\nu}{}_{b}V^{\mu}{}_{\nu}.
\tag{J.11}
\end{equation*}
回头看式 (J.5)，我们看到度规张量在正交归一基中的分量恰好就是平直度规的分量 $\eta_{ab}$。（正因如此，希腊指标有时被称为“弯曲的”（curved），拉丁指标则被称为“平直的”（flat）。）事实上，我们甚至可以用平直度规及其逆 $\eta^{ab}$ 来升降拉丁指标。你可以自行验证一切都行之有效（例如，用度规降低指标的操作与在正交归一基和坐标基之间切换是相互对易的）。特别地，我们对逆 vielbein 的定义与我们通常的指标升降规则是一致的：
\begin{equation*}
e^{\mu}{}_{a} = g^{\mu\nu}\eta_{ab}e_{\nu}{}^{b}.
\tag{J.12}
\end{equation*}

我们引入 vielbein $e_{\nu}{}^{a}$ 时，是把它们当作一组基矢量在另一个基下求出的分量。这等价于把它们看作一个 $(1,1)$ 型张量的分量：
\begin{equation*}
e = e_{\nu}{}^{a}\mathrm{d}x^{\nu}\otimes\hat{e}_{(a)}.
\tag{J.13}
\end{equation*}
但这其实是一个我们早已熟知且喜爱的张量：恒等映射。让这个张量作用在一个矢量上，得到的仍是同一个矢量，只不过换了一个基；这正是式 (J.10) 的内容。同样，如果用逆 vielbein $e^{\mu}{}_{a}$ 把 $e_{\nu}{}^{a}$ 上的拉丁指标换成希腊指标，根据式 (J.3) 就得到 Kronecker delta $\delta^{\mu}_{\nu}$，它当然就是作用在矢量（或 1-形式）上的恒等映射。这一点值得强调，因为我们也可以选择把 $e_{\nu}{}^{a}$ 解释为一组矢量分量（有些参考文献正是这么做的），那样的话协变导数的形式会有所不同。引入了一组新的基矢量和 1-形式之后，我们便不得不重提我们最爱讨论的话题：变换性质。一直以来我们都小心强调，张量变换律只是坐标变换的间接结果；真正的问题在于基的改变。如今我们有了非坐标基，这些基就可以独立于坐标而改变。唯一的限制是必须保持正交归一性条件 (J.1)。但我们知道什么样的变换保持平直度规——在欧几里得号差的度规下它们是正交变

% ---- 书页 486 / p501 ----
换，而在洛伦兹号差的度规下则是洛伦兹变换。因此我们考虑如下形式的基变换：
\begin{equation*}
\hat{e}_{(a)} \to \hat{e}_{(a')} = \Lambda^{a}{}_{a'}(x)\hat{e}_{(a)},
\tag{J.14}
\end{equation*}
其中矩阵 $\Lambda^{a}{}_{a'}(x)$ 表示依赖于位置的变换，它们（在每一点）保持度规的规范形式不变：
\begin{equation*}
\Lambda^{a}{}_{a'}\Lambda^{b}{}_{b'}\eta_{ab} = \eta_{a'b'}.
\tag{J.15}
\end{equation*}
事实上，这些矩阵对应于我们在平直空间中所说的逆洛伦兹变换（作用在基矢量上的那类）；与之前一样，我们也有普通的洛伦兹变换 $\Lambda^{a'}{}_{a}$，它变换基 1-形式。就分量而言，与之前一样，我们用 $\Lambda^{a'}{}_{a}$ 变换上指标，用 $\Lambda^{a}{}_{a'}$ 变换下指标。

于是，我们现在获得了在空间的每一点都实施一个洛伦兹变换的自由（视号差而定，也可以是普通的欧几里得转动）。这些变换因此被称为\textbf{局部洛伦兹变换}（local Lorentz transformations），简称 LLT。我们仍然保有通常的改变坐标的自由，这被称为\textbf{广义坐标变换}（general coordinate transformations），简称 GCT。两者可以同时发生，由此得到一个混合的张量变换律：
\begin{equation*}
T^{a'\mu'}{}_{b'\nu'} = \Lambda^{a'}{}_{a}\frac{\partial x^{\mu'}}{\partial x^{\mu}}\Lambda^{b}{}_{b'}\frac{\partial x^{\nu}}{\partial x^{\nu'}}T^{a\mu}{}_{b\nu}.
\tag{J.16}
\end{equation*}
把我们关于张量的知识翻译成非坐标基的语言，在多数情况下只不过是把 vielbein 塞到正确的位置而已。关键的例外出现在我们要对东西求导的时候。在我们通常的形式体系中，张量的协变导数等于它的偏导数再加上若干修正项——每个指标一项——涉及该张量与联络系数。同样的程序对非坐标基依然成立，只是要把通常的联络系数 $\Gamma^{\lambda}_{\mu\nu}$ 换成\textbf{自旋联络}（spin connection），记作 $\omega_{\mu}{}^{a}{}_{b}$。每个拉丁指标都以通常的方式配上一项自旋联络：
\begin{equation*}
\nabla_{\mu}X^{a}{}_{b} = \partial_{\mu}X^{a}{}_{b} + \omega_{\mu}{}^{a}{}_{c}X^{c}{}_{b} - \omega_{\mu}{}^{c}{}_{b}X^{a}{}_{c}.
\tag{J.17}
\end{equation*}
（“自旋联络”这个名字源于这样一个事实：可以用它来对旋量求协变导数，而用常规的联络系数这实际上是做不到的。）当拉丁指标和希腊指标混合出现时，我们会同时得到两种类型的项。

张量应当与其写法无关这一通常要求，让我们得以推导自旋联络、vielbein 与 $\Gamma^{\lambda}_{\mu\nu}$ 之间的关系。考虑矢量 $X$ 的协变导数，先在纯坐标基中：
\begin{align*}
\nabla X &= (\nabla_{\mu}X^{\nu})\mathrm{d}x^{\mu}\otimes\partial_{\nu} \\
&= (\partial_{\mu}X^{\nu} + \Gamma^{\nu}_{\mu\lambda}X^{\lambda})\mathrm{d}x^{\mu}\otimes\partial_{\nu}. \tag{J.18}
\end{align*}

% ---- 书页 487 / p502 ----
现在在混合基中求同一个对象，再转换到坐标基：
\begin{align*}
\nabla X &= (\nabla_{\mu}X^{a})\mathrm{d}x^{\mu}\otimes\hat{e}_{(a)} \\
&= (\partial_{\mu}X^{a} + \omega_{\mu}{}^{a}{}_{b}X^{b})\mathrm{d}x^{\mu}\otimes\hat{e}_{(a)} \\
&= \big(\partial_{\mu}(e_{\nu}{}^{a}X^{\nu}) + \omega_{\mu}{}^{a}{}_{b}e_{\lambda}{}^{b}X^{\lambda}\big)\mathrm{d}x^{\mu}\otimes(e^{\sigma}{}_{a}\partial_{\sigma}) \\
&= e^{\sigma}{}_{a}\big(e_{\nu}{}^{a}\partial_{\mu}X^{\nu} + X^{\nu}\partial_{\mu}e_{\nu}{}^{a} + \omega_{\mu}{}^{a}{}_{b}e_{\lambda}{}^{b}X^{\lambda}\big)\mathrm{d}x^{\mu}\otimes\partial_{\sigma} \\
&= \big(\partial_{\mu}X^{\nu} + e^{\nu}{}_{a}\partial_{\mu}e_{\lambda}{}^{a}X^{\lambda} + e^{\nu}{}_{a}e_{\lambda}{}^{b}\omega_{\mu}{}^{a}{}_{b}X^{\lambda}\big)\mathrm{d}x^{\mu}\otimes\partial_{\nu}. \tag{J.19}
\end{align*}
与式 (J.18) 比较可得
\begin{equation*}
\Gamma^{\nu}_{\mu\lambda} = e^{\nu}{}_{a}\partial_{\mu}e_{\lambda}{}^{a} + e^{\nu}{}_{a}e_{\lambda}{}^{b}\omega_{\mu}{}^{a}{}_{b},
\tag{J.20}
\end{equation*}
或者等价地
\begin{equation*}
\omega_{\mu}{}^{a}{}_{b} = e_{\nu}{}^{a}e^{\lambda}{}_{b}\Gamma^{\nu}_{\mu\lambda} - e^{\lambda}{}_{b}\partial_{\mu}e_{\lambda}{}^{a}.
\tag{J.21}
\end{equation*}
稍作处理，就可以把这个关系写成 vielbein 的协变导数为零：
\begin{align*}
\nabla_{\mu}e_{\nu}{}^{a} &= \partial_{\mu}e_{\nu}{}^{a} - \Gamma^{\lambda}_{\mu\nu}e_{\lambda}{}^{a} + \omega_{\mu}{}^{a}{}_{b}e_{\nu}{}^{b} \\
&= 0, \tag{J.22}
\end{align*}
它有时被称为“tetrad 假设”（tetrad postulate）。注意，它总是成立的；推导它并不需要对联络作任何假设。具体地说，我们不需要假设联络是度规相容的或无挠的。不过，我们确实隐含地取 $e_{\nu}{}^{a}$ 代表 $(1,1)$ 型张量 (J.13)；既然这个张量就是恒等映射，它的协变导数为零也就不足为奇了。（并非所有文献都持这种观点，所以要当心。）

既然联络可以被看作我们为修补协变导数的变换律而不得不引入的东西，那么自旋联络本身不服从张量变换律也就不足为怪了。实际上，在 GCT 之下，那个下希腊指标确实按正确的方式变换，即像 1-形式那样。但在 LLT 之下，自旋联络非齐次地变换：
\begin{equation*}
\omega_{\mu}{}^{a'}{}_{b'} = \Lambda^{a'}{}_{a}\Lambda^{b}{}_{b'}\omega_{\mu}{}^{a}{}_{b} - \Lambda^{c}{}_{b'}\partial_{\mu}\Lambda^{a'}{}_{c}.
\tag{J.23}
\end{equation*}
建议你自行验证，这确实能使协变导数按正确的方式变换。

到目前为止，我们做的不过是空洞的形式主义——把已经知道的东西翻译成一套新记号。但我们做的这些工作确实换来了两样东西。第一样我们已经提到过：能够描述时空上的旋量场并对它们求协变导数；这里我们不再深入。第二样是视角的改变：我们可以把各种张量看作张量值的微分形式（tensor-valued differential forms）。例如，像 $X_{\mu}{}^{a}$ 这样的对象——我们

% ---- 书页 488 / p503 ----
把它看作带混合指标的 $(1,1)$ 型张量——也可以看作一个“矢量值的 1-形式”（vector-valued one-form）。它有一个下希腊指标，所以我们把它当作 1-形式，但对下指标的每一个取值它又是一个矢量。类似地，张量 $A_{\mu\nu}{}^{a}{}_{b}$（在 $\mu$ 和 $\nu$ 上反对称）可以看作一个“(1,1) 型张量值的 2-形式”。于是，任何带有若干个反对称下希腊指标和若干个拉丁指标的张量，都可以看作一个微分形式，只是取值于张量丛。（普通的微分形式简单来说就是标量值的形式。）当我们考虑外微分时，这一视角的用处就显现出来。如果我们想把 $X_{\mu}{}^{a}$ 当作矢量值的 1-形式，我们自然会想取它的外微分：
\begin{equation*}
(\mathrm{d}X)_{\mu\nu}{}^{a} = \partial_{\mu}X_{\nu}{}^{a} - \partial_{\nu}X_{\mu}{}^{a}.
\tag{J.24}
\end{equation*}
容易验证，这个对象在 GCT 之下像 2-形式一样变换［也就是说，按照 $(0,2)$ 型张量的变换律］，但在 LLT 之下却不像矢量那样变换（洛伦兹变换依赖于位置，这会在变换律中引入一个非齐次项）。不过，我们可以明智地运用自旋联络来修复这一点；由于非张量性的变换律 (J.23)，自旋联络可以被看作一个 1-形式，却不是张量值的 1-形式。于是，对象
\begin{equation*}
(\mathrm{d}X)_{\mu\nu}{}^{a} + (\omega\wedge X)_{\mu\nu}{}^{a} = \partial_{\mu}X_{\nu}{}^{a} - \partial_{\nu}X_{\mu}{}^{a} + \omega_{\mu}{}^{a}{}_{b}X_{\nu}{}^{b} - \omega_{\nu}{}^{a}{}_{b}X_{\mu}{}^{b},
\tag{J.25}
\end{equation*}
你可以自行验证，它会像一个正规的张量那样变换。

这套形式体系的一个直接应用是挠率和曲率的表达式——它们是刻画任何给定联络的两个张量。挠率带有两个反对称下指标，可以看作矢量值的 2-形式 $T_{\mu\nu}{}^{a}$。曲率总是在最后两个指标上反对称，是 $(1,1)$ 型张量值的 2-形式 $R^{a}{}_{b\mu\nu}$。利用我们在微分形式上省略指标的自由，可以用基 1-形式把它们表达为
\begin{equation*}
e^{a} = e_{\mu}{}^{a}\mathrm{d}x^{\mu}
\tag{J.26}
\end{equation*}
以及自旋联络 1-形式
\begin{equation*}
\omega^{a}{}_{b} = \omega_{\mu}{}^{a}{}_{b}\mathrm{d}x^{\mu}.
\tag{J.27}
\end{equation*}
注意我们已经更换了记号，定义 $e^{a} \equiv \hat{\theta}^{(a)}$。这是相当惯用的做法，而且更简洁。于是挠率和曲率的定义关系就是
\begin{equation*}
T^{a} = \mathrm{d}e^{a} + \omega^{a}{}_{b}\wedge e^{b}
\tag{J.28}
\end{equation*}
以及
\begin{equation*}
R^{a}{}_{b} = \mathrm{d}\omega^{a}{}_{b} + \omega^{a}{}_{c}\wedge\omega^{c}{}_{b}.
\tag{J.29}
\end{equation*}

% ---- 书页 489 / p504 ----
请记住，$R^{a}{}_{b}$ 代表整个 Riemann 张量（希腊指标被省略了）；不要把它与 Ricci 张量混淆。这两个式子就是所谓的 \textbf{Cartan 结构方程}（Cartan structure equations）。它们与通常的定义等价；让我们就挠率作一番证明练习，曲率的验证留给你自己完成。我们有
\begin{align*}
T_{\mu\nu}{}^{\lambda} &= e^{\lambda}{}_{a}T_{\mu\nu}{}^{a} \\
&= e^{\lambda}{}_{a}\big(\partial_{\mu}e_{\nu}{}^{a} - \partial_{\nu}e_{\mu}{}^{a} + \omega_{\mu}{}^{a}{}_{b}e_{\nu}{}^{b} - \omega_{\nu}{}^{a}{}_{b}e_{\mu}{}^{b}\big) \\
&= \Gamma^{\lambda}_{\mu\nu} - \Gamma^{\lambda}_{\nu\mu}, \tag{J.30}
\end{align*}
这正是我们最初给出的定义。这里我们用到了式 (J.20)，即用 vielbein 和自旋联络表出的 $\Gamma^{\lambda}_{\mu\nu}$。我们还可以把这些张量所满足的恒等式表达为
\begin{equation*}
\mathrm{d}T^{a} + \omega^{a}{}_{b}\wedge T^{b} = R^{a}{}_{b}\wedge e^{b}
\tag{J.31}
\end{equation*}
以及
\begin{equation*}
\mathrm{d}R^{a}{}_{b} + \omega^{a}{}_{c}\wedge R^{c}{}_{b} - R^{a}{}_{c}\wedge\omega^{c}{}_{b} = 0.
\tag{J.32}
\end{equation*}
其中第一式是 $R^{\rho}{}_{[\sigma\mu\nu]} = 0$ 的推广，第二式则是 Bianchi 恒等式 $\nabla_{[\lambda}R^{\rho}{}_{|\sigma|\mu\nu]} = 0$。（有时这两个方程都被称作 Bianchi 恒等式。）

这些表达式的形式会诱使人几乎无法抗拒地去定义一个“协变外微分”（covariant-exterior derivative）：它作用在张量值形式上，先取通常的外微分，再添上带自旋联络的适当项——每个拉丁指标一项。虽然这里我们不打算这样做，但向这种诱惑屈服也无妨；事实上，式 (J.28) 的右边以及式 (J.31) 和式 (J.32) 的左边，正可以看作这样的协变外微分。但要小心，式 (J.29) 却不能这样看；你不能对自旋联络取任何一种协变导数，因为它不是张量。

到目前为止，我们的方程对一般联络都成立；现在看看对 Christoffel 联络会得到什么。无挠要求就是 (J.28) 为零；这并不会直接给出关于自旋联络系数的任何简单表述。度规相容性则表现为度规的协变导数为零：$\nabla g = 0$。当我们把度规用正交归一基表达——此时它的分量就是简单的 $\eta_{ab}$——就能看到这会导出什么：
\begin{align*}
\nabla_{\mu}\eta_{ab} &= \partial_{\mu}\eta_{ab} - \omega_{\mu}{}^{c}{}_{a}\eta_{cb} - \omega_{\mu}{}^{c}{}_{b}\eta_{ac} \\
&= -\omega_{\mu ab} - \omega_{\mu ba}. \tag{J.33}
\end{align*}
于是令它等于零就意味着
\begin{equation*}
\omega_{\mu ab} = -\omega_{\mu ba}.
\tag{J.34}
\end{equation*}

% ---- 书页 490 / p505 ----
这样，度规相容性就等价于自旋联络在其拉丁指标上的反对称性。（和之前一样，只有当两个指标都在上或都在下时，这种说法才有意义。）这两个条件合在一起，让我们可以把自旋联络用 vielbein 表达出来。这个解有一个显式公式，但实践中更容易的做法是直接求解无挠条件
\begin{equation*}
\omega^{a}{}_{b}\wedge e^{b} = -\mathrm{d}e^{a},
\tag{J.35}
\end{equation*}
再利用自旋联络的非对称性求出各个分量。

思考非坐标基的最好理由之一是：在某些情形——包括计算曲率张量——它们确实能带来巨大的简化。让我们通过一个简单例子看看这是如何运作的：一个空间平直的膨胀宇宙，其度规为
\begin{equation*}
ds^{2} = -\mathrm{d}t^{2} + a^{2}(t)\delta_{ij}\mathrm{d}x^{i}\mathrm{d}x^{j}.
\tag{J.36}
\end{equation*}
我们将使用式 (J.26) 和式 (J.27) 的微分形式记号；像这样的计算就是很好的证据，说明这套语言不仅优雅而且确实实用。于是度规可以写成（对任何几何都成立）
\begin{equation*}
ds^{2} = \eta_{ab}e^{a}\otimes e^{b}.
\tag{J.37}
\end{equation*}
我们需要选取基 1-形式 $e^{a}$，使得上式与我们的度规 (J.36) 相吻合。有多种选择（它们由局部洛伦兹变换相联系），但其中有一个是显而易见的：
\begin{align*}
e^{0} &= \mathrm{d}t \\
e^{i} &= a\,\mathrm{d}x^{i}. \tag{J.38}
\end{align*}
现在我们想利用式 (J.35) 解出自旋联络。好消息是，基本上靠猜测就能做到。首先，适当地升降指标（用 $\eta^{ab}$ 和 $\eta_{ab}$），我们导出 $\omega_{ab}$ 反对称性的推论：
\begin{align*}
\omega^{0}{}_{0} &= 0 \\
\omega^{0}{}_{j} &= \omega^{j}{}_{0} \\
\omega^{i}{}_{j} &= -\omega^{j}{}_{i}. \tag{J.39}
\end{align*}
接下来计算式 (J.35) 的右边：
\begin{align*}
\mathrm{d}e^{0} &= 0 \\
\mathrm{d}e^{i} &= \mathrm{d}a\wedge\mathrm{d}x^{i} = \dot{a}\,\mathrm{d}t\wedge\mathrm{d}x^{i}, \tag{J.40}
\end{align*}

% ---- 书页 491 / p506 ----
然后再看左边：
\begin{align*}
\omega^{0}{}_{b}\wedge e^{b} &= \omega^{0}{}_{j}\wedge e^{j} = a\,\omega^{0}{}_{j}\wedge\mathrm{d}x^{j} \\
\omega^{i}{}_{b}\wedge e^{b} &= \omega^{i}{}_{0}\wedge e^{0} + \omega^{i}{}_{j}\wedge e^{j} = \omega^{i}{}_{0}\wedge\mathrm{d}t + \omega^{i}{}_{j}\wedge\mathrm{d}x^{j}. \tag{J.41}
\end{align*}
代入式 (J.35) 得到
\begin{align*}
\omega^{0}{}_{j}\wedge\mathrm{d}x^{j} &= 0 \\
\omega^{i}{}_{0}\wedge\mathrm{d}t + \omega^{i}{}_{j}\wedge\mathrm{d}x^{j} &= -\dot{a}\,\mathrm{d}t\wedge\mathrm{d}x^{i}. \tag{J.42}
\end{align*}
我们想从这些方程解出 $\omega^{a}{}_{b}$。很容易想到猜 $\omega^{0}{}_{j} = 0$；但这样一来，为了解第二个方程就需要 $\omega^{i}{}_{j} = -\dot{a}\delta^{i}_{j}\mathrm{d}t$，这与式 (J.39) 中的 $\omega^{i}{}_{j} = -\omega^{j}{}_{i}$ 不相容。不过，考虑到楔积的反对称性，我们可以令 $\omega^{0}{}_{j}$ 正比于 $\mathrm{d}x^{j}$ 来解第一个方程。确实，如果选取
\begin{equation*}
\omega^{0}{}_{j} = \dot{a}\,\mathrm{d}x^{j}, \qquad \omega^{i}{}_{0} = \dot{a}\,\mathrm{d}x^{i},
\tag{J.43}
\end{equation*}
就会发现，只要令
\begin{equation*}
\omega^{i}{}_{j} = 0.
\tag{J.44}
\end{equation*}
式 (J.42) 中的两个方程便都得到满足。

既然已经知道自旋联络，我们就能通过下式轻松得到曲率：
\begin{equation*}
R^{a}{}_{b} = \mathrm{d}\omega^{a}{}_{b} + \omega^{a}{}_{c}\wedge\omega^{c}{}_{b}.
\tag{J.45}
\end{equation*}
我们先计算自旋联络形式的外微分：
\begin{align*}
\mathrm{d}\omega^{i}{}_{0} &= \ddot{a}\,\mathrm{d}t\wedge\mathrm{d}x^{i} \\
\mathrm{d}\omega^{0}{}_{j} &= \ddot{a}\,\mathrm{d}t\wedge\mathrm{d}x^{j} \\
\mathrm{d}\omega^{i}{}_{j} &= 0, \tag{J.46}
\end{align*}
然后计算楔积：
\begin{align*}
\omega^{0}{}_{c}\wedge\omega^{c}{}_{0} &= 0 \\
\omega^{i}{}_{c}\wedge\omega^{c}{}_{0} &= 0 \\
\omega^{i}{}_{c}\wedge\omega^{c}{}_{j} &= \dot{a}^{2}\mathrm{d}x^{i}\wedge\mathrm{d}x^{j}. \tag{J.47}
\end{align*}
于是我们得到曲率 2-形式：
\begin{align*}
R^{0}{}_{0} &= 0 \\
R^{0}{}_{j} &= \ddot{a}\,\mathrm{d}t\wedge\mathrm{d}x^{j} \\
R^{i}{}_{0} &= \ddot{a}\,\mathrm{d}t\wedge\mathrm{d}x^{i} \\
R^{i}{}_{j} &= \dot{a}^{2}\mathrm{d}x^{i}\wedge\mathrm{d}x^{j}. \tag{J.48}
\end{align*}

% ---- 书页 492 / p507 ----
为了便于比较，我们可以用 vielbein 把 $R^{a}{}_{b\mu\nu}$ 换算成我们惯用的表达式 $R^{\rho}{}_{\sigma\mu\nu}$：
\begin{equation*}
R^{\rho}{}_{\sigma\mu\nu} = e^{\rho}{}_{a}e_{\sigma}{}^{b}R^{a}{}_{b\mu\nu}.
\tag{J.49}
\end{equation*}
用分量形式写出，vielbein (J.38) 及其逆为
\begin{equation*}
e_{\mu}{}^{a} = \begin{pmatrix} 1 & & & \\ & a & & \\ & & a & \\ & & & a \end{pmatrix}, \qquad e^{\nu}{}_{b} = \begin{pmatrix} 1 & & & \\ & a^{-1} & & \\ & & a^{-1} & \\ & & & a^{-1} \end{pmatrix}. \tag{J.50}
\end{equation*}
我们还需要计算基形式的楔积的分量，这相当直接：
\begin{equation*}
(\mathrm{d}x^{\alpha}\wedge\mathrm{d}x^{\beta})_{\mu\nu} = \delta^{\alpha}_{\mu}\delta^{\beta}_{\nu} - \delta^{\alpha}_{\nu}\delta^{\beta}_{\mu}.
\tag{J.51}
\end{equation*}
把所有这些放在一起，就得到 $R^{\rho}{}_{\sigma\mu\nu}$ 的分量：
\begin{align*}
R^{0}{}_{j0l} &= a\ddot{a}\delta_{jl} \\
R^{i}{}_{0k0} &= -\frac{\ddot{a}}{a}\delta^{i}_{k} \\
R^{i}{}_{jkl} &= \dot{a}^{2}(\delta^{i}_{k}\delta_{jl} - \delta^{i}_{l}\delta_{jk}), \tag{J.52}
\end{align*}
以及在最后两个指标上作反对称得到的那些分量。我们可以缩并，得到 Ricci 张量 $R_{\sigma\nu} = R^{\lambda}{}_{\sigma\lambda\nu}$ 的分量：
\begin{align*}
R_{00} &= -3\frac{\ddot{a}}{a} \\
R_{i0} &= 0 \\
R_{ij} &= (a\ddot{a} + 2\dot{a}^{2})\delta_{ij}. \tag{J.53}
\end{align*}
你可以验证，这与我们在第 8 章得到的结果一致。仅在这个简单的例子里，tetrad 方法在计算上就已经比坐标基方法简单；对于更复杂的度规，这种相对优势还会继续扩大。

借助于非坐标基的语言，我们可以把黎曼几何中联络与曲率的形式体系，同粒子物理中规范理论的形式体系加以比较。在这两种情形下，我们感兴趣的场都生活在指派给时空每一点的矢量空间里。在黎曼几何中，这些矢量空间包括切空间、余切空间以及由它们构造的高阶张量空间。而在规范理论中，我们关心的是“内部”矢量空间。区别在于：切空间及其同类与流形本身紧密相连，一旦流形被搭建起来，它们便自然地得到定义；例如，切空间可以看作一点处方向导数的空间。与此相反，内部矢量

% ---- 书页 493 / p508 ----
空间则可以有我们想要的任何维数，并且必须作为流形之外的一种独立添加来定义。用数学行话说，底流形与内部矢量空间（在每一点定义）的并就是一个\textbf{纤维丛}（fiber bundle），而矢量空间的每一份拷贝都被称为一根“纤维”（fiber）（与我们定义切丛的方式一致）。

除了底流形（对我们而言就是时空）和纤维之外，纤维丛定义中的另一个重要要素是“结构群”（structure group）：一个作用在纤维上的李群，它描述各纤维在相互交叠的坐标片上如何被缝合在一起。不深入细节地说，四维时空中切丛的结构群一般是 $\mathrm{GL}(4,\mathbf{R})$，即实可逆 $4\times4$ 矩阵群；如果有洛伦兹型度规，它可以约化为洛伦兹群 $SO(3,1)$。现在设想我们引入一个内部三维矢量空间，并用普通的转动把纤维缝合起来；这个新丛的结构群便是 $SO(3)$。生活在该丛中的场可以记作 $\phi^{A}(x^{\mu})$，其中 $A$ 从一取到三；它在流形的每一点上都是一个三维矢量（内部的，与时空无关）。我们可以随意选择纤维中的基；这意味着“物理量”应当在如下形式的局部 $SO(3)$ 变换下保持不变：
\begin{equation*}
\phi^{A}(x^{\mu}) \to \phi^{A'}(x^{\mu}) = O^{A'}{}_{A}(x^{\mu})\phi^{A}(x^{\mu}),
\tag{J.54}
\end{equation*}
其中 $O^{A'}{}_{A}(x^{\mu})$ 是 $SO(3)$ 中一个依赖于时空的矩阵。这类变换被称为\textbf{规范变换}（gauge transformations），在其下不变的理论被称为“规范理论”（gauge theories）。

在多数情况下，不难把事情安排得让物理量在规范变换下不变。唯一的困难出现在我们考虑偏导数 $\partial_{\mu}\phi^{A}$ 的时候。由于矩阵 $O^{A'}{}_{A}(x^{\mu})$ 依赖于时空，它会给偏导数的变换贡献一个不想要的项。到现在你应该已经能猜到解决方案了：引入一个联络，来修正变换律中的非齐次项。于是，我们把纤维丛上的联络定义为对象 $A_{\mu}{}^{A}{}_{B}$：它带两个“群指标”和一个时空指标。在 GCT 之下它像 1-形式那样变换，而在规范变换之下它的变换律是
\begin{equation*}
A_{\mu}{}^{A'}{}_{B'} = O^{A'}{}_{A}O^{B}{}_{B'}A_{\mu}{}^{A}{}_{B} - O^{C}{}_{B'}\partial_{\mu}O^{A'}{}_{C}.
\tag{J.55}
\end{equation*}
（当心：我们的约定与粒子物理文献中的约定不同。）在这个变换律下，“规范协变导数”（gauge covariant derivative）
\begin{equation*}
D_{\mu}\phi^{A} = \partial_{\mu}\phi^{A} + A_{\mu}{}^{A}{}_{B}\phi^{B}
\tag{J.56}
\end{equation*}
在规范变换下“像张量那样”变换，欢迎你自行验证。［在普通的电磁学中，联络就是惯用的矢势。不需要任何指标，因为结构群 $U(1)$ 是一维的。］

% ---- 书页 494 / p509 ----
显然，内部纤维丛上联络的这一概念，与切丛上的联络密切相关，尤其是在我们一直讨论的正交归一标架图像中。例如，变换律 (J.55) 就与自旋联络的变换律 (J.23) 完全相同。我们还可以定义一个曲率张量，或称“场强”（field strength）张量，它是一个 2-形式：
\begin{equation*}
F^{A}{}_{B} = \mathrm{d}A^{A}{}_{B} + A^{A}{}_{C}\wedge A^{C}{}_{B},
\tag{J.57}
\end{equation*}
与式 (J.29) 完全对应。我们可以沿路径平行移动物体，也有一个与平行移动传播子类似的结构；把一个矢量沿闭合曲线平行移动一周，所得矩阵的迹被称为“Wilson 圈”（Wilson loop）。

我们本可以继续发展切丛与内部矢量丛之间的关系，但那就得另写一本书了。让我们转而通过强调这两种构造之间的重要\emph{差异}来收尾。差异源于这样一个事实：切丛与底流形紧密相关，而其他纤维丛则是事后才附加上去的。说 $p$ 点切空间中的一个矢量“指向沿一条穿过 $p$ 的路径的方向”是有意义的；但对内部矢量丛来说，这种说法毫无意义。因此，内部空间没有坐标基的类似物——沿曲线的偏导数与内部矢量毫无关系。依次推下去，也就没有类似 vielbein 那样把正交归一基与坐标基联系起来的东西。特别是挠率张量，只对切丛上的联络才有定义，对任何规范理论联络都没有；它可以被看作 vielbein 的协变外微分，而内部丛上没有这样的构造。你应当领会联络这一概念的不同用法之间的联系，但同时也别忘乎所以。

\section{习题}

% 习题编号照原书
\textbf{1.}\quad 在式 (J.37) 中我们提过，度规在正交归一基中可以写成
\begin{equation*}
ds^{2} = \eta_{ab}e^{a}\otimes e^{b}.
\tag{J.58}
\end{equation*}
这怎么可能呢？如果度规的分量处处都是 $\eta_{ab}$，我们怎么可能知道几何是什么？

\textbf{2.}\quad 计算 Mixmaster 宇宙（Mixmaster universe）的联络 1-形式、曲率 2-形式，进而求出 Riemann 张量的分量。度规为
\begin{equation*}
ds^{2} = -\mathrm{d}t\otimes\mathrm{d}t + \alpha^{2}\sigma^{1}\otimes\sigma^{1} + \beta^{2}\sigma^{2}\otimes\sigma^{2} + \gamma^{2}\sigma^{3}\otimes\sigma^{3}.
\end{equation*}
其中 $\alpha$、$\beta$、$\gamma$ 只是 $t$ 的函数，而 1-形式 $\sigma^{i}$ 为
\begin{align*}
\sigma^{1} &= \cos\psi\,\mathrm{d}\theta + \sin\psi\sin\theta\,\mathrm{d}\phi \\
\sigma^{2} &= \sin\psi\,\mathrm{d}\theta - \cos\psi\sin\theta\,\mathrm{d}\phi \\
\sigma^{3} &= \mathrm{d}\psi + \cos\theta\,\mathrm{d}\phi.
\end{align*}
